% =====================================================================
%  Security Auditing for Python Cryptographic Code
%  Rounds 1-6 Training, Data Governance, and Evaluation Report (EN)
%  Build: latexmk -pdf round1-6_experiment_report_en.tex
% =====================================================================
\documentclass[11pt]{article}

\usepackage[T1]{fontenc}
\usepackage[utf8]{inputenc}
\usepackage{lmodern}
\usepackage{microtype}
\usepackage[a4paper, margin=2.4cm]{geometry}
\usepackage{amsmath, amssymb}
\usepackage{booktabs}
\usepackage{array}
\usepackage{tabularx}
\usepackage{multirow}
\usepackage{graphicx}
\usepackage{xcolor}
\usepackage{seqsplit}
\usepackage{caption}
\usepackage{enumitem}
\usepackage{fancyvrb}
\usepackage[hidelinks]{hyperref}

\newcolumntype{L}[1]{>{\raggedright\arraybackslash}p{#1}}
\newcolumntype{Y}{>{\raggedright\arraybackslash}X}

\captionsetup{font=small, labelfont=bf}
\setlength{\parskip}{0.42em}
\setlength{\emergencystretch}{3em}
\setlist[itemize]{leftmargin=1.4em}
\setlist[enumerate]{leftmargin=1.7em}

\newcommand{\code}[1]{\texttt{\detokenize{#1}}}
\newcommand{\hash}[1]{\texttt{\seqsplit{#1}}}
\newcommand{\breakable}[1]{\path{#1}}
\newcommand{\good}{\textcolor{green!50!black}{$\checkmark$}}
\newcommand{\bad}{\textcolor{red!70!black}{$\times$}}

\title{\bfseries Security Auditing for Python Cryptographic Code\\
\large Rounds 1--6 Training, Data Governance, and Evaluation Report}
\author{AI-Assisted Cryptographic Code Security Auditor}
\date{2026-10-05}

\begin{document}

\maketitle
\vspace{-1.2em}

\begin{abstract}
\noindent This report documents a research prototype for auditing cryptographic misuse in Python.
The system uses 13 in-house Semgrep CRYPTO rules as a static candidate generator and a QLoRA fine-tuned
Qwen3-4B as the semantic discriminator, covering two tasks: \textbf{detect} (decide at function level
whether the code is vulnerable, and emit a CWE and severity) and \textbf{triage} (given code plus a static
finding, emit Confirm/Reject, an explanation, and an optional patch). Rounds 1--4 established the two-task
baseline, exposed catastrophic forgetting from incremental continued training, recovered through full
fresh-base retraining, and tested synthetic negative-sample scaling. Round 5 produced a dataset QA failure
and a clear model regression. Round 6 folded real-data strengthening, rule fixtures, strong-model pre-review,
and browser-based human review into a single full retrain. The shipped R6 artifacts contain 4,122
train/validation/test records, of which 401 are new and 79 are real (30 unique verified cryptographic
positives); the frozen holdout is 19 detect and 36 triage records (only 4 Confirm). R6 recovers from the R5
detect collapse, reaching F1 0.588 on the frozen holdout and 0.671 on the design test split, but it fails the
precision, FPR, and 18-probe accuracy gates. This report therefore defines R6 as a \textbf{reproducible
research baseline}, not a production-grade auditor. It deliberately preserves the Round 6 candidate
harvesting, cost estimation, K3 binary review, browser human review, provenance tagging, dataset freezing,
and failure analysis for later review and for R7 design.
\end{abstract}

\tableofcontents

\section{Project Scope and Purpose}

\subsection{One-line summary}

This project studies a composite code-audit prototype: static analysis finds suspicious cryptographic API
usages with high recall; a small domain model decides whether a finding is real, classifies the CWE, explains
the reason, and proposes a fix direction; uncertain records go to a strong model or to browser-based human
review. The system ultimately emits structured JSON and a Markdown audit report.

The long-term product hypothesis is a ``security-audit copilot'': it does not replace a security engineer,
rather it reduces manual triage time and keeps the audit rationale, data provenance, and model verdicts as
traceable evidence.

\subsection{The two tasks}

\begin{table}[h]
\centering
\caption{The system's two audit tasks}
\label{tab:tasks-en}
\small
\begin{tabularx}{\linewidth}{L{2.2cm}Y Y}
\toprule
\textbf{Task} & \textbf{Input} & \textbf{Output} \\
\midrule
detect & one Python function or code snippet & \code{vulnerable}, \code{cwe}, \code{severity}, \code{confidence}, \code{explanation} \\
triage & code + a static-analysis finding & \code{cwe}, \code{severity}, \code{verdict=Confirm|Reject}, \code{explanation}, optional \code{patch} \\
\bottomrule
\end{tabularx}
\end{table}

\subsection{Scope and non-goals}

The current version covers only the defined cryptographic-misuse patterns: weak hashing, broken algorithms,
AES-ECB, predictable IV/nonce, hard-coded keys, weak randomness, weak key length, insecure TLS, plaintext
passwords, weak KDF parameters, and non-constant-time comparison. It does not claim to discover unknown
classes of cryptographic vulnerabilities, and it does not generate attacks or exploit code.

This report explicitly does not claim any of the following:

\begin{itemize}
  \item that the model reaches production-grade precision, recall, or CWE routing;
  \item that strong-model K3 review was rewritten into human ground truth;
  \item that the merged training artifacts are leak-free;
  \item that R6's test-split boundary values represent true cryptographic generalisation;
  \item that ``training longer'' is a solution to the current bottleneck.
\end{itemize}

\section{Background: Cryptography, CWE, and Code Auditing}

\subsection{Common cryptographic building blocks}

Auditing cryptographic code requires separating ``the algorithm itself'' from ``how the algorithm is used in
the business logic''. A weak primitive is not automatically a vulnerability: MD5 for password storage is
usually a defect, whereas MD5 for a cache key or change detection usually is not. This context boundary is
exactly what the model finds hardest to learn.

\begin{table}[h]
\centering
\caption{Cryptographic building blocks and their audit concerns}
\label{tab:crypto-basics-en}
\small
\begin{tabularx}{\linewidth}{L{2.7cm}Y Y}
\toprule
\textbf{Building block} & \textbf{Role} & \textbf{Common audit risks} \\
\midrule
Symmetric encryption & Encrypt/decrypt with one key, e.g.\ AES & ECB mode, fixed IV, nonce reuse, weak key length \\
AEAD & Confidentiality plus integrity, e.g.\ AES-GCM, ChaCha20-Poly1305 & Nonce reuse breaks authentication; wrong key origin and nonce lifetime \\
Asymmetric encryption and signatures & Public-key encryption, private-key signing, e.g.\ RSA, ECDSA, Ed25519 & Small RSA modulus, wrong padding, weak key or randomness source \\
Hashing & Map arbitrary input to a fixed-length digest, e.g.\ SHA-256 & MD5/SHA-1 for security-sensitive data; mistaking a plain digest for a password hash \\
MAC/HMAC & Keyed message integrity and authenticity & Non-constant-time comparison, key reuse, algorithm downgrade \\
KDF and password hashing & Derive keys or resist offline cracking, e.g.\ PBKDF2, scrypt, Argon2 & Too few iterations, missing salt, attacker-controlled parameters \\
Randomness & Generate keys, tokens, salts, IVs, nonces & Using \code{random} (Mersenne Twister) instead of \code{secrets} or \code{os.urandom} \\
TLS and certificates & Verify server identity and establish an encrypted channel & \code{verify=False}, \code{CERT_NONE}, hostname check disabled, trust root overridden \\
\bottomrule
\end{tabularx}
\end{table}

IV (initialization vector) and nonce (number used once) are prerequisites for the security of many modes.
CBC-like modes require an unpredictable IV; streaming AEAD such as GCM requires that a nonce never repeat
under the same key. Merely seeing ``the function uses AES'' is not enough to judge safety; the mode, IV/nonce,
key origin, and intended use must all be checked.

Password hashing must not be confused with general-purpose hashing. A bare \code{sha256(password)} is fast,
which is fine for file integrity but unsuitable for resisting offline cracking. A password system should use
a salted, cost-tunable KDF or password hash; modern choices include Argon2id, scrypt, bcrypt, or a
standard-conforming PBKDF2 configuration. The audit must also judge whether the iteration count is too low,
whether the salt is random and unique, and whether the parameters can be attacker-controlled.

\subsection{CWE taxonomy and the 13 CRYPTO rules}

The project binds static rules to CWE targets. A CWE is a defect class; a rule is an engineering-detectable
pattern; the two are not identical. One rule can cover several CWEs and one CWE may need several rules.

\begin{table}[h]
\centering
\caption{The 13 CRYPTO rule families and their core CWEs}
\label{tab:cwe13-en}
\small
\begin{tabularx}{\linewidth}{L{2.35cm}Y L{2.35cm}}
\toprule
\textbf{Rule} & \textbf{Problem} & \textbf{Target CWE} \\
\midrule
CRYPTO-001 & Weak digest (MD5 etc.) in a security-sensitive context & CWE-327 / CWE-916 \\
CRYPTO-002 & Weak digest (SHA-1 etc.) in a security-sensitive context & CWE-327 / CWE-916 \\
CRYPTO-003 & DES or 3DES & CWE-327 \\
CRYPTO-004 & RC4 & CWE-327 \\
CRYPTO-005 & AES-ECB & CWE-327 \\
CRYPTO-006 & Predictable or reused IV/nonce & CWE-329 \\
CRYPTO-007 & Hard-coded key or password & CWE-321 \\
CRYPTO-008 & Weak pseudo-random number generator & CWE-338 \\
CRYPTO-009 & Weak key length or weak curve & CWE-326 \\
CRYPTO-010 & Insecure TLS or certificate validation & CWE-295 \\
CRYPTO-011 & Plaintext password storage or logging & CWE-256 \\
CRYPTO-012 & Weak KDF parameters & CWE-916 \\
CRYPTO-013 & Non-constant-time comparison & CWE-208 \\
\bottomrule
\end{tabularx}
\end{table}

The report focuses on nine target CWEs: CWE-321 (hard-coded key), CWE-326 (inadequate encryption strength),
CWE-327 (broken or risky algorithm), CWE-329 (IV/nonce not randomized), CWE-338 (weak PRNG), CWE-916
(insufficient password-hash computational effort), CWE-208 (observable timing discrepancy), CWE-295
(improper certificate validation), and CWE-256 (plaintext storage of a password). The R6 candidate pool also
surfaces neighbouring classes such as CWE-798 (hard-coded credentials), CWE-330 (insufficient randomness),
and CWE-328 (reversible one-way hash), but neighbouring classes must not enter the target training set without
a rule mapping.

\subsection{SAST, Semgrep, AST, and false positives}

SAST (static application security testing) analyses source code without running the program. Semgrep produces
findings through pattern and lightweight semantic matching; the Python AST extracts functions, compares
structure, locates calls, and detects near-duplicates.

The strength of static analysis is speed, interpretability, and precise location; its weakness is that a
pattern hit is not the same as a real defect. For example:

\begin{itemize}
  \item a crypto library implements MD5 or DES for legacy-protocol or historical-format compatibility, which
        is a ``valid false positive'';
  \item test fixtures, cache keys, ETags, consistency checks, sharding, and jitter use weak primitives
        usually without a security objective;
  \item a function does call AES but uses a fixed IV, ECB, or a predictable nonce, which is still a real
        problem.
\end{itemize}

The system therefore adopts ``static recall plus LLM semantic routing'': the static rules act as a candidate
generator, and the LLM judges purpose and security boundary. The LLM can reduce false positives, but it can
also hallucinate cache, test, or ETag context that is absent from the code, so a structured schema, an
evaluation set, and provenance auditing are mandatory.

\subsection{How detect and triage differ}

detect is an open judgement: the model sees only the code, decides whether it is vulnerable, and gives a CWE.
triage is a constrained judgement: the model sees both the code and one specific finding, and decides Confirm
or Reject. Their supervision signals differ:

\begin{itemize}
  \item detect needs safe code, real vulnerabilities, and both rule-hit and rule-miss negatives;
  \item triage needs Reject cases (rule hit but no security purpose) and Confirm cases (rule hit with a real
        security-sensitive purpose);
  \item the detect training distribution controls whether the model ``over-reports everything'';
  \item the triage training distribution controls whether the model over-generalises weak-primitive usages
        into Reject.
\end{itemize}

\subsection{LLM fine-tuning basics}

This project uses supervised fine-tuning (SFT) and QLoRA. SFT treats the input prompt and the expected JSON
output as a training pair; LoRA freezes the base weights and trains only low-rank adapter matrices; QLoRA
further applies 4-bit quantisation to reduce memory. The core hyperparameters are:

\begin{table}[h]
\centering
\caption{Fine-tuning terminology}
\label{tab:finetune-en}
\small
\begin{tabularx}{\linewidth}{L{3.0cm}Y}
\toprule
\textbf{Term} & \textbf{Meaning and effect} \\
\midrule
rank (r) & Rank of the LoRA low-rank matrices; larger means more capacity and more risk of overfitting or disturbing old abilities \\
alpha & LoRA scaling coefficient; usually proportional to rank; R1--R6 mainly use rank=16, alpha=32 \\
dropout & Random-dropout rate of the adapter layers; used to suppress overfitting \\
learning rate & Update step size; too large destroys base abilities, too small learns no new boundary \\
epochs & Number of passes over the data; more is not better and can reinforce templates and noise \\
max length & Truncation length; over-long samples are truncated, losing code context \\
fresh-base retrain & Retrain from the original base with the full merged dataset; the most stable paradigm in this project \\
incremental update & Continue training the existing model on new samples only; R2 and R4 both triggered catastrophic forgetting \\
catastrophic forgetting & After learning a new task, performance on old tasks degrades significantly \\
\bottomrule
\end{tabularx}
\end{table}

Catastrophic forgetting recurred throughout the project. R2 fed only 54 triage increments, dropped probe FPR
to 0, but pushed detect CWE accuracy from 0.796 down to 0.433 and domain Confirm recall from 0.357 down to 0.
R4 micro-continued training was worse: detect recall fell from 0.990 to 0.584 and domain Confirm recall from
0.929 to 0.071. The engineering constraint that follows is: \textbf{when tasks or distributions change
noticeably, do a fresh-base full merge retrain; incremental continued training is only for small repairs on
the same distribution.}

\subsection{Evaluation metrics}

Let TP, FP, FN, and TN denote true positive, false positive, false negative, and true negative. The main
detect metrics are:

\begin{align}
\mathrm{precision} &= \frac{TP}{TP+FP},&
\mathrm{recall} &= \frac{TP}{TP+FN},\\
F1 &= \frac{2\cdot \mathrm{precision}\cdot \mathrm{recall}}{\mathrm{precision}+\mathrm{recall}},&
\mathrm{accuracy} &= \frac{TP+TN}{TP+FP+FN+TN}.
\end{align}

CWE accuracy is computed only over ground-truth positives and asks whether the CWE is exactly correct once
the code has been judged vulnerable. triage Confirm recall is the fraction of true Confirm cases retained;
FPR (false positive rate) is the fraction of true Reject cases that are wrongly Confirmed:

\begin{equation}
\mathrm{FPR}=\frac{FP}{FP+TN}.
\end{equation}

Small positive counts make recall extremely unstable. The frozen triage set has only 4 Confirm cases, so a
single prediction flip moves recall by 25 points; this report therefore treats frozen Confirm recall as a
count, not an acceptance gate. Likewise, 30 real cryptographic positives is below the project's pre-set
50--100 threshold, so aggregate scores cannot be used to claim real-world cryptographic generalisation.

\section{System Architecture, Data Sources, and Provenance Tagging}

\subsection{Audit pipeline}

\begin{center}
\fbox{\parbox{0.94\linewidth}{\small
Python source
$\rightarrow$ 13 Semgrep CRYPTO rules
$\rightarrow$ candidate finding
$\rightarrow$ Qwen3-4B fine-tuned model (detect / triage)
$\rightarrow$ structured JSON
$\rightarrow$ Markdown audit report
}}
\end{center}

The main modules are \code{rules/}, \code{analyzer/}, \code{agent/}, \code{model/}, \code{evaluation/},
\code{scripts/}, and \code{reviewer/}. Training and inference are governed by the same prompt templates, so
that the ChatML training data and the deployment input distribution stay consistent.

\subsection{Data layers}

\begin{table}[h]
\centering
\caption{Semantic layers of the training data}
\label{tab:data-layers-en}
\small
\begin{tabularx}{\linewidth}{L{2.5cm}L{2.2cm}Y}
\toprule
\textbf{Layer} & \textbf{Typical source} & \textbf{Use and limitation} \\
\midrule
Public vulnerability data & PyCode-Vul & Provides general Python detect signal; most entries are Web/application bugs, not crypto-specific truth \\
Domain rule fixtures & Semgrep, Bandit test cases & Provide rule hits, safe alternative forms, Reject/negatives; \code{verified=false}, not real vulnerability samples \\
Manually/model-constructed samples & curated, combination templates & Fill the 13 rule families and class ratios; may be templated, so overfitting must be prevented \\
Real advisory records & GHSA, OSV/PyPI, PyPA, fix commit & Vulnerability and fix function pairs with a provenance chain; currently the most valuable real crypto signal \\
Secure reference & Fixed-commit secure-library code & Used for fix suggestions and safe-code comparison; does not enter training unless it passes a review gate \\
\bottomrule
\end{tabularx}
\end{table}

\subsection{R6 final data snapshot}

R6 merges real records directly into the sixth round rather than creating a separate Round 7. The final
uploaded set is:

\begin{table}[h]
\centering
\caption{R6 final data split}
\label{tab:r6-split-en}
\small
\begin{tabular}{lrrr}
\toprule
\textbf{Split} & \textbf{detect} & \textbf{triage} & \textbf{Total} \\
\midrule
train & 2732 & 383 & 3115 \\
val & 299 & 48 & 347 \\
test & 602 & 58 & 660 \\
\midrule
\textbf{Total} & \textbf{3633} & \textbf{489} & \textbf{4122} \\
\bottomrule
\end{tabular}
\end{table}

R6 adds 401 records relative to the R5 base:

\begin{table}[h]
\centering
\caption{R6 added-record provenance}
\label{tab:r6-new-en}
\small
\begin{tabularx}{\linewidth}{L{4.0cm}rY}
\toprule
\textbf{Source layer} & \textbf{Count} & \textbf{verified and use} \\
\midrule
curated, rule fixtures, and supplements & 322 & All \code{verified=false}; used for rule coverage, negatives/Reject, and tool regression \\
Real advisory detect records & 30 & All vulnerability positives; 23 by K3, 7 by human \\
Real advisory triage records & 49 & 30 Confirm / 19 Reject; 37 by K3, 12 by human \\
\midrule
\textbf{Total} & \textbf{401} & 79 real records; unique real crypto positives still only 30 \\
\bottomrule
\end{tabularx}
\end{table}

Of the 79 real records, 60 come from K3 model review and 19 from browser human review. Model-reviewed records
keep \code{review_mode=model} and \code{model=kimi-k3}; human records keep \code{review_mode=manual}. Even
though the project permits K3 labels as final labels, the report must not describe them as human review.

The CWE distribution of real detect positives is: CWE-295 (12), CWE-208 (9), CWE-327 (3), CWE-338 (2),
CWE-321 (2), CWE-326 (1), CWE-256 (1). The real triage export is 30 Confirm / 19 Reject, with the signal
concentrated in CRYPTO-010 (12 Confirm) and CRYPTO-013 (9 Confirm).

\subsection{Data-composition risk: class balance is not effective signal}

R6 detect has 2,037 positives and 1,596 negatives; the surface ratio looks reasonable, but only 434 positives
belong to target crypto CWEs and the other 1,603 are general Web/application/memory vulnerabilities. In other
words, most positive supervision is spent on non-crypto defects. The per-CWE sample counts are:

\begin{table}[h]
\centering
\caption{Target crypto CWE counts among detect positives}
\label{tab:cwe-count-en}
\small
\begin{tabular}{lr@{\hspace{1.4cm}}lr}
\toprule
\textbf{CWE} & \textbf{Positives} & \textbf{CWE} & \textbf{Positives} \\
\midrule
CWE-327 & 275 & CWE-916 & 11 \\
CWE-326 & 63 & CWE-256 & 10 \\
CWE-295 & 33 & CWE-321 & 5 \\
CWE-208 & 18 & CWE-329 & 5 \\
CWE-338 & 14 & & \\
\bottomrule
\end{tabular}
\end{table}

Per-class recall on rare CWEs will have very high variance. If training samples randomly over the total, rare
classes and real advisory positives are diluted further; future work needs a train-only crypto-focus sampler
with minimum quotas by CWE and rule.

\subsection{Leakage audit}

The goal of the frozen evaluation set is to remove every sample whose whitespace-normalised code and AST are
identical to something in the final R5/R6 train or val set. The frozen R6 evaluation set is:

\begin{itemize}
  \item detect: 19 records (7 positive / 12 negative), 0 removed;
  \item triage: 36 records (4 Confirm / 32 Reject), with 6 records removed as duplicates of
        \code{round4-real-passlib}.
\end{itemize}

However, a cross-split audit of the merged artifacts still finds residual duplicates inherited from R4/R5:

\begin{table}[h]
\centering
\caption{Residual near-duplicates in the merged training artifacts}
\label{tab:leakage-en}
\small
\begin{tabular}{lrr}
\toprule
\textbf{Pair} & \textbf{Exact duplicate rows} & \textbf{Identical-AST rows} \\
\midrule
train / val & 2 & 21 \\
train / test & 2 & 17 \\
val / test & 2 & 2 \\
\bottomrule
\end{tabular}
\end{table}

It is therefore only defensible to claim that ``the frozen evaluation set is not exposed to the optimiser-visible
set'', not that the entire merged training set is leak-free. Future additions must keep splitting by repository
and commit, not randomly by row.

\section{Rounds 1--5: How the Problems Emerged}

\subsection{Round 1: two-task baseline}

\textbf{Goal}: first domain adaptation on public Python vulnerability data plus domain crypto data, running the
full detect/triage chain.

\textbf{Data}: detect 2,406 + triage 132 = 2,538 records.

\textbf{Key configuration}: LoRA rank=8, learning rate 1e-4.

\textbf{Results}:

\begin{itemize}
  \item detect CWE accuracy 0.796, F1 0.656;
  \item domain Confirm recall 0.357;
  \item 18 triage probes: Confirm recall 0.917, FPR 0.667, accuracy 0.722.
\end{itemize}

\textbf{Exposed problem}: the model was usable on general detect, but triage Confirmed many non-security uses
of weak primitives. Typical false positives included MD5 in cache keys, SHA-1 in change detection,
\code{random} in backoff jitter, test fixtures, and legacy-compatible implementations. FPR 0.667 shows that
``rule hit'' was wrongly treated as ``real security defect''.

\subsection{Round 2: triage incremental training and catastrophic forgetting}

\textbf{Goal}: add triage increments only, to push down FPR.

\textbf{Data}: 54 records (28 Confirm + 26 Reject), triage only.

\textbf{Key configuration}: batch=8, learning rate 1e-4, 3 epochs, rank=16 / alpha=32 / dropout=0.05,
max length=1024.

\textbf{Results}:

\begin{itemize}
  \item probe FPR fell from 0.667 to \textbf{0.000} and accuracy rose from 0.722 to 0.944;
  \item detect CWE accuracy fell from 0.796 to \textbf{0.433}, with 343/569 samples shifting CWE prediction;
  \item domain Confirm recall fell from 0.357 to \textbf{0.000}, with the model judging almost everything
        Reject.
\end{itemize}

\textbf{Lesson}: feeding only new-task increments destroys the old-task representation. A short-term
improvement in a target metric is not an overall capability gain. Incremental continued training was from then
on listed as a high-risk route.

\subsection{Round 3: fresh-base full merge retrain}

\textbf{Goal}: abandon R2 continued training, retrain from the base on the full merged dataset, and restore
both detect and triage.

\textbf{Data}: 2,595 = R1's 2,538 + R2's 54 + 3 signature/authentication repairs.

\textbf{Key configuration}: batch=16, learning rate 1e-4, 3 epochs, rank=16 / alpha=32 / dropout=0.005,
cosine, max length=2048.

\textbf{Key decision}: append \code{TRIAGE_PURPOSE_GUIDE} to the triage prompt, requiring the model to judge
explicitly whether the flagged call is for a security or a non-security purpose. Without the guide the same
model had FPR 0.500; with the guide FPR was 0.000 while Confirm recall stayed 1.000.

\textbf{Results}:

\begin{itemize}
  \item detect recall 0.990, precision 0.537, F1 0.697, CWE accuracy 0.921;
  \item domain Confirm recall 0.929, FPR 0.067, accuracy 0.931;
  \item 18 probes: accuracy 18/18.
\end{itemize}

\textbf{Leftover problem}: detect still reported almost everything, 562/569 judged vulnerable, with
\code{tp/fp/fn/tn = 302/260/3/4}. Precision 0.537 is nearly the positive prior 0.536, showing that the binary
vulnerable decision lacks discriminative power. The root cause is that detect negatives were almost all
non-crypto Django code, so the model had never seen ``secure cryptographic code''.

\subsection{Round 4: negative-sample scaling and real-project validation}

\textbf{Goal}: add secure cryptographic negatives, reduce detect over-reporting, and hold out real crypto
libraries for validation.

\textbf{Negative-sample design}:

\begin{itemize}
  \item Class A: modern secure patterns such as scrypt, Argon2, bcrypt, PBKDF2, AES-GCM with random nonce,
        ChaCha20-Poly1305, SHA-3, HKDF, ECDH, \code{secrets}/\code{os.urandom}, RSA-OAEP, Ed25519;
  \item Class B: weak primitives used for non-security purposes such as cache keys, deduplication, sharding,
        change detection, backoff jitter, and test fixtures;
  \item strong-B: rules do fire, but the model must judge non-vulnerable from the purpose;
  \item weak-B: rules do not fire; used as ordinary safe negatives.
\end{itemize}

\textbf{Two routes}:

\begin{table}[h]
\centering
\caption{R4 incremental continued training vs full retrain}
\label{tab:r4-routes-en}
\small
\begin{tabularx}{\linewidth}{L{2.4cm}Y Y}
\toprule
 & \textbf{Micro continued training} & \textbf{Full fresh-base} \\
\midrule
Start point & R3 model & original base \\
Data & 77 negatives & 3,315 (detect 3126 + triage 189) \\
Negatives & 42 A + 35 B & added 150 positives + 493 negatives, final 1563:1563 \\
lr / epochs & 5e-5 / 2 & 1e-4 / 3 \\
\bottomrule
\end{tabularx}
\end{table}

\textbf{Results}:

\begin{itemize}
  \item full 50:50: detect recall 0.984, precision 0.536, F1 0.694, CWE accuracy 0.897; fp=260, tn=4 basically
        unchanged from R3;
  \item micro continued training: detect recall 0.990 $\rightarrow$ 0.584, domain Confirm recall 0.929
        $\rightarrow$ 0.071, catastrophic forgetting again;
  \item 42 real-project validation: detect judged only 4/7 safe functions correctly; domain triage FPR was
        0.429 (micro) / 0.571 (full); all 11 passlib legacy-hash cases were wrongly judged Confirm.
\end{itemize}

\textbf{Lesson}: the number of synthetic negatives did not teach the model the security boundary of real
libraries. Templated negatives are easily memorised and cannot replace real library code and
protocol-compatibility scenarios.

\subsection{Round 5: dataset QA failure and model regression}

R5 was meant to keep repairing the data distribution and triage ability, but the round produced a clear
regression. The submitted R5 run on the detect slice produced \code{tp/fp/fn/tn = 0/2/3/8}, F1 0.000,
accuracy 0.615; triage n=41, 3 Confirm / 38 Reject, FPR 0.368. The R5 dataset QA report is FAIL and includes:

\begin{itemize}
  \item \code{backtest_pos_missing=22}: samples that should have a positive back-test signal are missing;
  \item \code{backtest_bneg_no_fire=25}: strong-B negatives did not fire the expected rule;
  \item the same configuration re-run produced detect \code{0/1/3/9} and triage FPR 0.289, showing run-to-run
        variance in the point estimate.
\end{itemize}

\textbf{Lesson}: the data-quality gate must precede training. R5's regression was part data problem and part
run variance; a single metric must not be treated as stable capability, and bad data must not enter the next
round unfixed.

\section{Round 6: Data Optimisation, Strong-Model Review, and Human Review}

\subsection{R6 repair goals}

R6 was designed against the four R5 problems:

\begin{table}[h]
\centering
\caption{R5 problems and R6 countermeasures}
\label{tab:r6-fix-en}
\small
\begin{tabularx}{\linewidth}{L{4.1cm}Y Y}
\toprule
\textbf{R5 problem} & \textbf{R6 countermeasure} & \textbf{Result} \\
\midrule
Triage Confirm recall decline & Add Confirm samples, keep binary review & R6 adds triage Confirm 113 / Reject 72 \\
Domain FPR rebound & Scale Reject and hedge with secure libraries & Frozen FPR 0.125, still above the 0.06 target \\
Triage class imbalance & Add samples by rule family and type & All 13 rule families covered \\
Insufficient detect rule coverage & Add rule fixtures and real advisory records & 13/13 rule families covered; 30 real positives still low \\
\bottomrule
\end{tabularx}
\end{table}

\subsection{The 322 new curated/rule-fixture records}

The 322 baseline additions come from 8 fixed source files:

\begin{table}[h]
\centering
\caption{Sources of the 322 R6 baseline additions}
\label{tab:r6-322-en}
\small
\begin{tabular}{lr}
\toprule
\textbf{Source file} & \textbf{Count} \\
\midrule
\code{new_detect.jsonl} & 33 \\
\code{new_triage.jsonl} & 21 \\
\code{supplement_detect.jsonl} & 28 \\
\code{supplement_triage.jsonl} & 5 \\
\code{supplement_v2_detect.jsonl} & 21 \\
\code{supplement_v2_triage.jsonl} & 6 \\
\code{rule_fixture_detect.jsonl} & 104 \\
\code{rule_fixture_triage.jsonl} & 104 \\
\midrule
\textbf{Total} & \textbf{322} \\
\bottomrule
\end{tabular}
\end{table}

Detect is 216 records (130 positive / 86 negative) and triage is 185 records (113 Confirm / 72 Reject) once
real records are included; within the 322 baseline, detect is 186 (100 pos / 86 neg) and triage is 136
(83 Confirm / 53 Reject).

The rule fixtures mainly come from:

\begin{itemize}
  \item Semgrep official rule tests, using \code{\# ruleid:} / \code{\# ok:} markers; licensed under the
        Semgrep Rules License v1.0;
  \item Bandit example programs, licensed under Apache-2.0;
  \item only rules semantically consistent with CRYPTO-001 through CRYPTO-013 are mapped; broader patterns
        such as blowfish, RC2, IDEA, XOR, MD2, MD4, JWT, and sha224 are skipped;
  \item positives must fire the corresponding rule scope, and negatives must clearly use a safe alternative
        such as SHA-256, AES-GCM, \code{os.urandom}, or \code{create_default_context}.
\end{itemize}

These 208 rule fixtures are all \code{verified=false}. They are tool-regression and finding-coverage data,
not real vulnerability positives; redistribution must respect the licence of the Semgrep source lines.

\subsection{Secure-library hedging and the reference layer}

A hard problem in R4/R5 was that secure library code was heavily over-reported. R6 adds negatives from many
secure libraries, including cryptography, PyNaCl, itsdangerous, PyJWT, PyCryptodome, passlib, argon2-cffi,
ecdsa, google-auth, and oscrypto. The goal is not ``the project is secure, so the function is safe'' but a
per-function check of parameters, purpose, key origin, and context.

Of the 15 fixed-commit audit candidates in \code{result.txt}, 7 were marked \code{secure_reference} and
written to a separate reference layer; they keep \code{verified=false} and \code{training_eligible=false} and
are used only for fix suggestions, safe-code comparison, and later manual promotion, not for training. The
other 8 were recommended for exclusion and were not merged.

\subsection{The real candidate pool: from advisory to function-level code pair}

R6's data engineering is not ``download a dataset'' but a traceable pipeline of candidate building,
screening, code completion, review, and merging:

\begin{enumerate}
  \item collect candidates from GHSA, OSV/PyPI, PyPA, PyCode-Vul, and local JSON/JSONL;
  \item screen by CWE, rule, package name, commit, repo, and keyword; de-duplicate into a candidate pool;
  \item fetch vulnerable/fixed code pairs via the fix commit; the 366 that pass \code{ast.parse} are better
        suited to function-level training;
  \item rank by evidence strength into P1/P2/P3; P1/P2 enter the priority review queue;
  \item run strict binary pre-review with a strong model, or send to browser-based human review;
  \item only Confirm/Reject conclusions enter \breakable{real_verified_detect.jsonl} and
        \breakable{real_verified_triage.jsonl};
  \item merge the R5 base with the builder, producing train/val/test and the SHA-256 manifest.
\end{enumerate}

The key offline snapshot:

\begin{table}[h]
\centering
\caption{R6 candidate-pool snapshot}
\label{tab:candidate-snapshot-en}
\small
\begin{tabularx}{\linewidth}{L{6.2cm}Y}
\toprule
\textbf{Metric} & \textbf{Value} \\
\midrule
On-disk external advisory sources & GHSA approx.\ 380k JSON, OSV-PyPI approx.\ 26k JSON, PyPA 7,721 YAML, Bandit, Semgrep, Trail of Bits rules \\
Per-source harvester output & GHSA 534, OSV-PyPI 843, PyPA 706, PyCode-Vul 389 \\
Merge and de-duplicate & 2,718 merged, 1,861 selected after de-duplication \\
P1 / P2 / P3 & 5 / 2 / 1,854 \\
With code pair & 400 records, of which 366 parse at both ends with \code{ast.parse} \\
Missing code only & 549 records have a fix commit \\
Missing commit and repo & 703 records need a package-to-repo mapping first \\
Class B training volume & 331 parseable real production code pairs, no advisory \\
Full P3 strong-model estimate & K3 high approx.\ CNY 975.18, K2.7-code high approx.\ CNY 267.01; not affordable to run blind on everything \\
\bottomrule
\end{tabularx}
\end{table}

This snapshot is itself an important fact: R6's bottleneck is not ``no advisories'' but the gap between an
advisory and a trainable function pair, in code, commit, repository, and human evidence.

\subsection{Strong-model pre-review: budget, model choice, and strict binary output}

Brute-force full review is not realistic, so the project estimates cost before calling. The actual K3 usage for
the Class A function-candidate batch was:

\begin{table}[h]
\centering
\caption{K3 real-candidate review cost}
\label{tab:k3-cost-en}
\small
\begin{tabular}{lr}
\toprule
\textbf{Item} & \textbf{Value} \\
\midrule
Candidates & 60 \\
prompt tokens & 240,546 \\
completion tokens & 26,995 \\
Actual input cost & CNY 4.81092 \\
Actual output cost & CNY 2.69950 \\
\textbf{Actual total cost} & \textbf{CNY 7.51042} \\
Model / reasoning effort & \code{kimi-k3} / \code{high} \\
Price basis & input CNY 20 per million tokens, output CNY 100 per million tokens \\
\bottomrule
\end{tabular}
\end{table}

The estimation stage used a conservative assumption of about CNY 33.59 for high-reasoning output; the actual
output was far below the assumed ceiling. This gap shows why ``estimate first, call second, report actual usage
third'' is a necessary flow; cost must not be guessed from candidate count.

The review prompt uses a strict binary format:

\begin{itemize}
  \item only \code{confirmed} or \code{rejected} is allowed;
  \item insufficient evidence is forced to \code{rejected}; no uncertain option is introduced;
  \item API, parsing, or schema errors are recorded as error and not imported;
  \item imported records keep \code{review_mode=model}, \code{model=kimi-k3}, \code{model_confidence}, and
        \code{source_audit};
  \item the 60 K3 results must not be described as human ground truth.
\end{itemize}

The import report shows: 60 review lines; on incremental import 48 were already reviewed and 1 was newly
imported; 8 lacked vulnerable code and 3 had a vulnerable snippet that is not valid Python; all were skipped
and never entered the verified export. The final K3 contribution is 23 detect + 37 triage.

\subsection{The browser-based human review platform}

To reduce manual effort, only the uncertain and sampled parts go to browser review. The final UI is
deliberately simplified to two buttons:

\begin{itemize}
  \item \textbf{Confirm}: exports one detect positive and one triage Confirm;
  \item \textbf{Reject}: exports one triage Reject;
  \item uncertain, scope-out, and license fields are removed;
  \item Confirm requires a target rule, target CWE, an explanation, and two Python snippets that pass
        \code{ast.parse};
  \item Reject requires a vulnerable snippet and an explanation;
  \item the reviewer may fix the snippet in the browser, but cannot enter samples unsupported by the source.
\end{itemize}

Each verdict is the source of truth, stored at:

\begin{Verbatim}[fontsize=\small]
data/round6/review/web/<reviewer>/labels/<verdict>/<candidate_id>.json
\end{Verbatim}

The application rebuilds \breakable{real_verified_detect.jsonl} and
\breakable{real_verified_triage.jsonl} from these
files. The final 19 human-reviewed records are 7 detect + 12 triage; one candidate was excluded because its
code exceeded the allowed length. Human labels and model labels share the same export schema but carry
different provenance fields.

\subsection{Pre-training freeze and reproducibility constraints}

Before training, R6 generates \code{r6_manifest.json}, freezing the source files and the SHA-256 of
train/val/test. The key artifacts are:

\begin{table}[h]
\centering
\caption{R6 uploaded-artifact hashes}
\label{tab:r6-hash-en}
\small
\begin{tabularx}{\linewidth}{L{2.4cm}rY}
\toprule
\textbf{Artifact} & \textbf{Rows} & \textbf{SHA-256} \\
\midrule
train & 3115 & \hash{139966e5217d4b2976db7a13f86d70ecb8aad09adc1198092f782e028041bd4d} \\
val & 347 & \hash{6b6797b5fb1da66c797f2efe28dc7559d124b2cd4adb5f75496a40f07d397f4b} \\
test & 660 & \hash{1549c1d7d5c8484aafe677ddf83ccfbdf6912e6052d1e910d49f45ed66c7d64c} \\
\bottomrule
\end{tabularx}
\end{table}

After freezing, train/val/test must not be modified. Evaluation uses \code{validation_frozen/}, which is
invisible during training.

\section{R6 Training and Reproducibility Information}

\subsection{Model and platform}

\begin{table}[h]
\centering
\caption{R6 model and workspace}
\label{tab:r6-model-en}
\small
\begin{tabularx}{\linewidth}{L{4.1cm}Y}
\toprule
\textbf{Item} & \textbf{Value} \\
\midrule
Deployed model ID & \code{qwen3-4b-instruct-2507-5f8261ad123d} \\
Fine-tune job & \code{ft-202610051556-cc87} \\
Fine-tune output model & \code{qwen3-4b-instruct-2507-ft-202610051556-cc87} \\
Base & \code{qwen3-4b-instruct-2507}, fresh base \\
Bailian workspace & \code{ws-avxkjb2tq5lq1gwm} \\
Training platform & Alibaba Cloud Bailian cn-beijing \\
Inference path & detect direct; triage appends \code{TRIAGE_PURPOSE_GUIDE} \\
\bottomrule
\end{tabularx}
\end{table}

\subsection{Hyperparameters and configuration differences}

The guide written before R6 recorded planned parameters, while the platform finally reported actual runtime
parameters. The two agree on LoRA structure but differ on batch and max length:

\begin{table}[h]
\centering
\caption{R6 planned hyperparameters vs platform report}
\label{tab:r6-hparams-en}
\small
\begin{tabular}{lcc}
\toprule
\textbf{Parameter} & \textbf{R6 guide plan} & \textbf{Platform / eval-report value} \\
\midrule
LoRA rank & 16 & 16 \\
LoRA alpha & 32 & 32 \\
LoRA dropout & 0.005 & 0.005 \\
learning rate & 1e-4 & 1e-4 \\
epochs & 3 & 3 \\
batch size & 16 & 128 \\
max length & 2048 & 4096 \\
scheduler & cosine & cosine \\
\bottomrule
\end{tabular}
\end{table}

This report uses the platform-reported parameters in the final model description while keeping the planned
values as a reproducibility difference. The difference also reminds us that an experiment report must record
both ``planned configuration'' and ``actual runtime configuration'', not only one of them.

\section{R6 Evaluation Results}

\subsection{Evaluation suites}

R6 uses four evaluation suites:

\begin{enumerate}
  \item frozen holdout: de-leaked 19 detect + 36 triage;
  \item 18 real triage probes: 12 Confirm + 6 Reject;
  \item R6 design test split: 602 detect + 58 triage, same source as training, valid only as a relative
        reference;
  \item capability matrix and secure-library anti-false-positive: 13 rules with one positive and one
        negative each, plus 10 secure patterns.
\end{enumerate}

\subsection{Frozen holdout}

\begin{table}[h]
\centering
\caption{R6 frozen holdout results}
\label{tab:r6-frozen-en}
\small
\begin{tabularx}{\linewidth}{l c c c c c Y}
\toprule
\textbf{Slice} & \textbf{n} & \textbf{recall} & \textbf{precision} & \textbf{F1} & \textbf{CWE acc} & \textbf{Note} \\
\midrule
detect & 19 & 0.714 & 0.500 & 0.588 & 0.400 & tp/fp/fn/tn = 5/5/2/7 \\
triage & 36 & 1.000$^{\ast}$ & -- & -- & -- & FPR 0.125, accuracy 0.889, only 4 Confirm \\
\bottomrule
\end{tabularx}
\smallskip
{\footnotesize $^{\ast}$ Confirm recall. With only 4 positives this is a count, not a statistical estimate.}
\end{table}

Of the 12 safe detect samples, 5 were over-reported; of the 7 vulnerability positives, 2 were missed. The 4
triage false positives were \code{Crypto_IO__PBES__PBES2_encrypt}, \code{passlib_cisco_pix},
\code{passlib_sun_md5_crypt}, and \code{r5v-011-tneg}. All 4 true Confirm cases were hit, but because there
are only 4 positives, Confirm recall 1.000 is only a count and cannot serve as a statistical acceptance test.

\subsection{The 18 real triage probes}

\begin{table}[h]
\centering
\caption{R6 18-probe results}
\label{tab:r6-probe-en}
\small
\begin{tabular}{lr}
\toprule
\textbf{Metric} & \textbf{Value} \\
\midrule
n & 18 (12 Confirm / 6 Reject) \\
Confirm recall & 0.500 \\
FPR & 0.000 \\
accuracy & 0.667 \\
\bottomrule
\end{tabular}
\end{table}

All six errors are Confirm $\rightarrow$ Reject; no Reject was over-reported as Confirm. The errors fall into
two classes:

\begin{enumerate}
  \item \textbf{explanation-verdict contradiction (3 cases)}: the explanation concludes \code{-> Confirm}
        but the JSON \code{verdict} outputs Reject, for example MD5 password comparison, a weak-random token,
        and an MD5 signature;
  \item \textbf{hallucinated absent security context (3 cases)}: the model explains a pure cryptographic
        function as a ``cache key'', ``test fixture generating an ETag'', or ``cache consistency'', while the
        input code has none of that context.
\end{enumerate}

This shows that R6's most serious regression is not a plain recall miss but an unstable triage decision chain:
explanation ability is still present, but the output field and the security-purpose reasoning are not fully
aligned.

\subsection{R6 design test split}

\begin{table}[h]
\centering
\caption{R6 test-split results}
\label{tab:r6-testsplit-en}
\small
\begin{tabularx}{\linewidth}{L{2.2cm}r r r Y}
\toprule
\textbf{Slice} & \textbf{n} & \textbf{recall} & \textbf{precision} & \textbf{F1 / FPR} \\
\midrule
detect & 602 (325 pos / 277 neg) & 0.859 & 0.550 & F1 0.671 \\
triage & 58 (28 Confirm / 30 Reject) & 0.571$^{\ast}$ & -- & FPR 0.200, accuracy 0.690 \\
\bottomrule
\end{tabularx}
\smallskip
{\footnotesize $^{\ast}$ Confirm recall.}
\end{table}

detect is \code{tp/fp/fn/tn = 279/228/46/49}, with CWE accuracy 0.724. The test split shares a source with
training, and the detect ground truth contains PyCode-Vul model-prediction noise, so this suite can only serve
as a relative reference.

\subsection{13-rule capability matrix}

\begin{table}[h]
\centering
\caption{R6 13-rule capability matrix}
\label{tab:r6-capability-en}
\small
\begin{tabularx}{\linewidth}{L{2.1cm}L{2.0cm}cL{1.9cm}c}
\toprule
\textbf{Rule} & \textbf{Expected CWE} & \textbf{Vuln.\ detected} & \textbf{Predicted CWE} & \textbf{Safe not flagged} \\
\midrule
CRYPTO-001 MD5 & CWE-327 & \good & CWE-327 & \good \\
CRYPTO-002 SHA-1 & CWE-327 & \bad & -- & \good \\
CRYPTO-003 DES & CWE-327 & \good & CWE-327 & \good \\
CRYPTO-004 RC4 & CWE-327 & \good & CWE-327 & \good \\
CRYPTO-005 AES-ECB & CWE-327 & \good & CWE-327 & \good \\
CRYPTO-006 predictable IV & CWE-329 & \good & CWE-327 & \good \\
CRYPTO-007 hard-coded key & CWE-321 & \good & CWE-327 & \bad \\
CRYPTO-008 weak randomness & CWE-338 & \bad & -- & \good \\
CRYPTO-009 weak key length & CWE-326 & \good & CWE-326 & \good \\
CRYPTO-010 insecure TLS & CWE-326 & \good & CWE-326 & \good \\
CRYPTO-011 plaintext password & CWE-256 & \good & CWE-259 & \bad \\
CRYPTO-012 weak KDF parameters & CWE-916 & \good & CWE-798 & \bad \\
CRYPTO-013 timing comparison & CWE-208 & \good & CWE-259 & \good \\
\bottomrule
\end{tabularx}
\smallskip
{\footnotesize The capability fixture labels CRYPTO-010 with expected CWE-326, although the rule table maps it
to CWE-295; this fixture/spec mismatch is itself a data-consistency item to fix in R7.}
\end{table}

Vulnerability detection is 11/13, CWE correctness is 6/13, and safe-not-flagged is 10/13. The model catches
obvious patterns but collapses several CWEs into CWE-327, CWE-259, or CWE-798, and still over-reports safe
patterns.

\subsection{Secure-library anti-false-positive}

Of 10 secure patterns, 8 were not over-reported, FPR 0.20. The false positives are:

\begin{itemize}
  \item \code{cryptography.hazmat...AESGCM(key).encrypt(os.urandom(12), ...)} judged vulnerable CWE-327;
  \item \code{nacl.secret.SecretBox(key).encrypt(msg)} judged vulnerable CWE-327.
\end{itemize}

\subsection{Acceptance gates}

\begin{table}[h]
\centering
\caption{R6 acceptance gates}
\label{tab:r6-gates-en}
\small
\begin{tabularx}{\linewidth}{L{3.7cm}Y L{2.2cm}}
\toprule
\textbf{Gate} & \textbf{Target} & \textbf{R6 status} \\
\midrule
detect precision & $>0.54$ & \bad\ frozen 0.500; test split 0.550, borderline \\
domain FPR & $\leq 0.06$ & \bad\ frozen 0.125, test split 0.200, secure libraries 0.20 \\
triage Confirm recall & report only & frozen has only 4 Confirm, no statistical meaning \\
18-probe accuracy & $\geq 0.95$ & \bad\ 0.667 \\
real-repository FPR & no worse than R3/R4 & \bad\ unresolved \\
\bottomrule
\end{tabularx}
\end{table}

\section{R6 Versus R5: What Recovered, What Did Not}

\subsection{Recovery}

\begin{table}[h]
\centering
\caption{Directional R5 vs R6 comparison}
\label{tab:r5-r6-en}
\small
\begin{tabularx}{\linewidth}{L{3.2cm}L{2.6cm}L{2.9cm}Y}
\toprule
\textbf{Metric} & \textbf{R5} & \textbf{R6} & \textbf{Interpretation} \\
\midrule
detect F1 & 0.000 & 0.588 frozen / 0.671 test & recovered from the detect collapse \\
detect recall & 0.000 & 0.714 frozen / 0.859 test & basic detection ability is back \\
triage FPR & 0.368 & 0.125 frozen & frozen-set false positives fell but remain above target \\
triage accuracy & 0.658 & 0.889 frozen & improved on the frozen set, but the composition differs \\
Confirm recall & 1.000 on 3 positives & 0.500 on 18 probes & shows small-positive recall cannot be extrapolated \\
\bottomrule
\end{tabularx}
\end{table}

This comparison is directional, not strictly matched: the R5 and R6 evaluation slices are not identical, and
R5 itself had run-to-run variance. R6 demonstrates that a full fresh-base retrain can recover from failure and
learn a weak signal for obvious cryptographic patterns, but it does not demonstrate production-grade precision
or real-world Confirm recall.

\subsection{Remaining problems}

\begin{enumerate}
  \item \textbf{CWE routing is a clear weakness}: frozen CWE accuracy 0.400, matrix 6/13, with a tendency to
        collapse into CWE-327, CWE-259, and CWE-798;
  \item \textbf{triage Confirm recall regressed}: 0.500 on the 18 probes, 0.571 on the test split, with
        explanation-verdict contradictions;
  \item \textbf{secure libraries are still over-reported}: AESGCM and \code{nacl.SecretBox} were wrongly
        flagged;
  \item \textbf{too few real positives}: 30 is below the 50--100 gate and cannot support a real-world crypto
        generalisation claim;
  \item \textbf{insufficient frozen-set statistical power}: detect 19 and triage 36 with only 4 Confirm;
  \item \textbf{prompt-alignment problem}: training data strips \code{TRIAGE_PURPOSE_GUIDE} while inference
        appends it; the guide ends with ``non-security purpose $\rightarrow$ Reject'', which may pull the
        verdict toward Reject;
  \item \textbf{training-distribution generalisation problem}: too many added Reject, bneg, legacy, test
        fixture, and cache samples; the model over-generalises these contexts onto real vulnerabilities;
  \item \textbf{residual leakage}: the frozen set is clean relative to the optimiser-visible set, but the
        merged train/val/test still contains inherited duplicates.
\end{enumerate}

\section{Conclusion, Limitations, and the R7 Roadmap}

\subsection{Conclusion}

R6 is an honest, reproducible research baseline. It shows that:

\begin{itemize}
  \item static rules + a small domain model + structured output can run end to end;
  \item a fresh-base full merge retrain can recover the R5 collapse;
  \item the model can learn explicit patterns such as MD5, DES, RC4, ECB, and weak TLS;
  \item data provenance, review mode, SHA-256 freezing, and de-leaked evaluation can be engineered;
  \item real-data bottleneck, sample volume, and contextual labels are constraints more critical than
        ``training longer''.
\end{itemize}

It does not show that:

\begin{itemize}
  \item production-grade precision, Confirm recall, or CWE routing is achieved;
  \item unknown cryptographic vulnerabilities can be discovered by generalisation;
  \item false positives on secure library code are robustly controlled;
  \item model review equals human ground truth;
  \item the merged training data is leak-free.
\end{itemize}

\subsection{Main limitations}

First, there are only 30 real crypto positives, concentrated in CWE-295 and CWE-208. Per-class results on rare
CWEs are unstable and cannot be hidden behind aggregate scores.

Second, the frozen set is very small. triage has only 4 Confirm cases, so a single prediction flip changes
recall noticeably; absolute counts and confidence intervals should be reported, and 1.000 must not be treated
as proof of capability.

Third, general PyCode-Vul labels are mixed with crypto-specific truth. Of 2,037 detect positives only 434
belong to target crypto CWEs, so the aggregate detect metric cannot measure crypto generalisation.

Fourth, strong-model review is model evidence, not human evidence. The 60 K3 records and 19 human records must
be counted separately; K3 can reduce workload but cannot replace independent ground truth.

Fifth, train/inference prompt inconsistency directly affects the verdict. R6's three ``correct explanation but
opposite verdict'' cases are likely related to this, and later work must either align the prompt or add a
structured field-consistency check.

Sixth, residual AST duplicates show that data inheritance still carries risk. Future work must split
stratified by repo/commit and re-run the exact/AST audit after every merge.

\subsection{R7 roadmap}

\begin{enumerate}
  \item \textbf{Scale real positives to 50--100}: continue GHSA/OSV/PyPA $\rightarrow$ fix commit $\rightarrow$
        parent source $\rightarrow$ function rebuild $\rightarrow$ strong-model/human review; report coverage
        by advisory, repo, CWE, and rule, not just total count.
  \item \textbf{Build a true external test set}: use time splits, repository splits, and fixed commits not
        involved in rule design; report bootstrap confidence intervals rather than a single point score.
  \item \textbf{Add a train-only crypto-focus sampler}: set minimum sampling quotas by target CWE and rule
        family, protect rare classes such as CWE-321/329/916/256/208, and never touch val/test.
  \item \textbf{Align train and inference prompts}: either add the purpose guide during training or remove it
        at inference; also add JSON-schema validation and an ``explanation consistent with verdict'' check.
  \item \textbf{Improve secure negative quality}: keep extracting from fixed-commit secure libraries, but
        check the call context first; prioritise AESGCM, SecretBox, legacy-format, and test/cache/ETag
        contexts so that weak-primitive usage is not over-generalised.
  \item \textbf{Improve CWE targets}: build per-CWE precision/recall for each target CWE and add CWE hierarchy
        or rule-mapping assistance so the model does not freely collapse to CWE-327.
  \item \textbf{Separate synthetic, rule-fixture, and real-record evaluation layers}: rule fixtures for tool
        regression, real records for quality evaluation, synthetic samples for coverage; the three must not
        share one headline.
  \item \textbf{Upgrade the base model or continue pre-training}: only after data quality passes the gate
        should a stronger base, a longer context, or domain continued pre-training be considered; a base-model
        upgrade cannot compensate for wrong labels.
  \item \textbf{Build an annotation quality mechanism}: double-review a small sample and measure agreement;
        perform error analysis on K3-versus-human disagreements; gradually convert high-confidence model
        labels into sampled labels.
  \item \textbf{Continue leakage and licence governance}: Semgrep rule fixtures follow the Semgrep Rules
        License v1.0 and Bandit follows Apache-2.0; retain provenance and licence metadata on publication.
\end{enumerate}

\section{Review Q\&A and Interview Points}

\subsection{Frequently asked questions}

\textbf{Why cryptographic code security?}

Cryptographic code has clear APIs and clear CWEs, and it is strongly context-dependent. It tests both static
rules and the LLM's weaknesses in security semantics, hallucination, and data quality.

\textbf{Why combine Semgrep and an LLM?}

Semgrep provides high-recall, locatable, interpretable candidates; the LLM provides purpose judgement and
false-positive filtering. Semgrep alone produces valid false positives from legacy compatibility, caches, and
test fixtures; an LLM alone lacks stable location and reproducible rules.

\textbf{Why fine-tune a small model instead of using a large-model API directly?}

A small model is cheaper, deployable, and researchable for domain adaptation, and it is easy to combine with
static rules. The project puts the large model in the data-review stage, not in the production inference path.

\textbf{What are the most important technical findings?}

First, fresh-base full retraining is more stable than incremental continued training; second, the data-quality
gate matters more than sample count; third, real crypto positives and human/model provenance tagging are the
core bottleneck; fourth, prompt alignment directly affects the triage verdict; fifth, small positive counts
make recall estimates very unstable.

\textbf{Was R6 a failure or a success?}

R6 successfully recovered the R5 detect collapse and formed a reproducible data and evaluation pipeline, but
it did not pass the precision, FPR, and probe-accuracy gates, so it can only be defined as a temporary research
baseline, not a successfully shipped production model.

\subsection{Things to watch during the experiments}

\begin{itemize}
  \item never mix metrics across rounds; every number must state its round, slice, and sample size;
  \item always report the positive count alongside Confirm recall;
  \item never write \code{verified=false} rule fixtures as real vulnerabilities;
  \item never write K3 labels as human labels;
  \item never claim zero leakage;
  \item never treat the R6 test-split borderline precision as evidence of production readiness;
  \item verify the manifest hashes before training and do not modify train/val/test afterwards;
  \item store both planned and actual hyperparameter values;
  \item check regulation, licence, and provenance for every real-data addition.
\end{itemize}

\appendix

\section{Reproduction Commands and Key Files}

\subsection{Build and freeze}

\begin{Verbatim}[fontsize=\small]
python3 scripts/build_round5_final_dataset.py --emit
python3 scripts/harvest_real_crypto_candidates.py \
  --advisories data/external/advisory-database/advisories
python3 scripts/screen_real_crypto_candidates.py --emit
python3 scripts/rebuild_candidate_functions.py --dry-run
GITHUB_TOKEN=ghp_xxx python3 scripts/rebuild_candidate_functions.py
python3 scripts/audit_real_crypto_with_llm.py --estimate-only
python3 scripts/annotate_real_crypto.py \
  --candidates data/round6/candidates/real_crypto_candidates_uncertain.jsonl \
  --reviewer reviewer-kimi --open
python3 scripts/build_round6_final_dataset.py --emit
python3 scripts/audit_dataset_leakage.py --write-frozen
python3 scripts/build_dataset_metadata.py
\end{Verbatim}

\subsection{Evaluation}

\begin{Verbatim}[fontsize=\small]
python3 scripts/run_round6_eval.py
python3 scripts/run_round6_test_split.py
python3 scripts/run_round6_capability.py
\end{Verbatim}

\subsection{Key artifacts}

\begin{itemize}
  \item \code{data/round6/eval/RESULTS.md}: full R6 evaluation conclusion;
  \item \code{data/round6/eval/benchmark_r6.json}: frozen detect + triage per-record predictions;
  \item \code{data/round6/eval/triage_probe_results_r6.json}: the 18 probes and explanations;
  \item \code{data/round6/eval/test_split_r6.json}: per-record predictions over the 660-record test split;
  \item \code{data/round6/eval/capability_r6.json}: the 13-rule matrix and secure-library probe;
  \item \code{reports/data_quality/model_audit_k3_cost.json}: real K3 cost;
  \item \code{reports/data_quality/model_audit_k3_import.json}: K3 import and skip reasons;
  \item \code{data/round6/final/r6_manifest.json}: source files and artifact hashes;
  \item \code{data/round6/validation_frozen/manifest.json}: frozen evaluation set and de-leak notes;
  \item \code{reviewer/README.md}: browser binary-review instructions.
\end{itemize}

\section{Glossary}

\begin{table}[h]
\centering
\small
\begin{tabularx}{\linewidth}{L{3.4cm}Y}
\toprule
\textbf{Term} & \textbf{Meaning} \\
\midrule
SAST & Static application security testing; analyse source without running the program \\
detect & Input code only; judge vulnerability and CWE \\
triage & Input code and a finding; decide Confirm/Reject \\
CWE & Common Weakness Enumeration \\
CSPRNG & Cryptographically secure pseudo-random number generator \\
AEAD & Authenticated encryption providing both confidentiality and integrity \\
KDF & Key-derivation function from a password or master key \\
LoRA / QLoRA & Low-rank adaptation / quantised low-rank adaptation \\
fresh-base retrain & Retrain from the original base on the full merged dataset \\
catastrophic forgetting & New training degrades old abilities \\
FPR & Fraction of true negatives wrongly judged positive \\
verified=false & Not yet independently verified for authenticity \\
\code{review_mode} & Review provenance: \code{manual} or \code{model} \\
\bottomrule
\end{tabularx}
\end{table}

\end{document}
